% !TEX TS-program = xelatex
% !TEX encoding = UTF-8 Unicode
% !BIB TS-program = bibtex

% 可积结构作为可学习的归纳偏置 —— 中文版（初稿）
% 命令行编译：./build.sh 或 latexmk -xelatex main.tex
% 中文支持依赖 ctex 包（TeX Live / MacTeX 默认携带）。

\documentclass[11pt]{article}

% --- 中文支持 -------------------------------------------------------------
\usepackage[UTF8, scheme=plain, fontset=fandol]{ctex}

% --- 数学 / 字体 ---------------------------------------------------------
\usepackage{amsmath, amssymb, amsfonts, amsthm}
\usepackage{mathrsfs}

% --- 引用 / 参考文献 -----------------------------------------------------
\usepackage[numbers]{natbib}

% --- 表格 / 图 -----------------------------------------------------------
\usepackage{array}
\usepackage{booktabs}
\usepackage{tabularx}
\usepackage{graphicx}
\usepackage{float}
\usepackage{subcaption}

% --- 排版 ----------------------------------------------------------------
\usepackage[margin=0.9in]{geometry}
\usepackage{verbatim}

% --- 超链接（最后载入）---------------------------------------------------
\usepackage[colorlinks=true, citecolor=blue, linkcolor=blue, urlcolor=blue,
            hypertexnames=false]{hyperref}

% --- 定理环境 ------------------------------------------------------------
\newtheorem{Theorem}{定理}[section]
\newtheorem{Definition}[Theorem]{定义}
\newtheorem{Proposition}[Theorem]{命题}
\newtheorem{Lemma}[Theorem]{引理}
\newtheorem{Remark}[Theorem]{备注}
% 实证结果（数值发现，非数学定理）用直立体，避免把实验包装成定理。
\theoremstyle{definition}
\newtheorem{Result}{实证结果}
\theoremstyle{plain}

% --- 元数据 --------------------------------------------------------------
% TODO（定稿）：最终题目待定。当前工作题目。
\title{可积结构作为可学习的归纳偏置：\\
       物理模型库的一次端到端样板}
\author{张栋\thanks{电子邮件：dongzhanghz@gmail.com}}
\date{\today}

\begin{document}
\maketitle

% --- 本文自定义记号 ------------------------------------------------------
\newcommand{\amp}{\mathrm{amp}}

% ============ 写作护栏（开工前务必读，定稿前删掉本注释块）============
%  护栏1（边界 / related work 必须处理，否则被一句话打回）：
%    - “循环/状态追踪能长度外推、注意力不能”是成熟结论（Liu 2023；Merrill 2023/2024）。
%      本文不以“赢 Transformer”为卖点。
%    - 物理归纳偏置有先例（HNN Greydanus 2019；LNN Cranmer 2020；可逆网 NICE/RealNVP）。
%      明说 delta：这些多为“把物理当目标/给保证”；本文测的是“可积结构让网络*学*出别的
%      模型+外挂都学不到的相互作用内容”，且用 leak/bolt-on 两个审计守诚实。
%    - 箱球系统是成熟数学物理对象（Takahashi–Satsuma 1990）；无人拿它当可学习序列模型比。
%  护栏2（范围，不可越界）：
%    - 主张只限“离散可积”。Toda（连续）是诚实负结果，写在 §边界，用来划界，不藏。
%    - 卖点是“精确性 / 相互作用内容”，不是“守恒”——守恒可外挂（bolt-on 对照证明）。
%    - 不说“建成了完整 benchmark”：现在是 validated Tier 1（单系统、K=3）。
%  护栏3（诚实标注，逐条）：数值硬验证（多种子） / 数值示意 / 展望。两个审计（leak、bolt-on）
%    是全文的诚实方法论，每张主表都要带。
% ==================================================================

\begin{abstract}
\noindent
% TODO：定稿前再压一遍。当前为诚实定位版。
AI 架构的诞生长期带有强烈的经验搜索色彩：模型常在局部改造、跨任务迁移和 benchmark 试错中
被筛出，而不是从一个系统的候选结构库中被选择。本文把物理学视为这样一座候选模型库，并用
\textbf{可积系统}这一格给出一次端到端样板：从物理结构出发，做成可训练网络，设计能检验其结构
优势的任务，并用审计对照排除硬编与外挂约束。具体地，我们在\textbf{离散}可积系统上给出一个
正面结果：在\emph{有限容量箱球系统}上，一个把\textbf{箱球的可积 carrier 结构}（带结构守恒的
进位扫描 $+$ 可学习放球门）做进去的模型，学会在多次碰撞的长组合下\textbf{精确保持孤子振幅内容}
（一个比球总数更丰富的守恒不变量）；\emph{在本文测试的模型、训练预算与评测协议下}，通用序列
模型（GRU / LSTM / Transformer）做不到，把守恒量\emph{外挂}到通用模型上也做不到。我们用两个
审计守住诚实：\textbf{泄露审计}
（把可学习的门换成写死规则）证明这一优势是\emph{真学出来的}而非硬编（写死规则只到 $40\%$
振幅内容，学出来的到 $100\%$）；\textbf{外挂对照}证明它\emph{外挂不出来}（钉住球总数
$100\%$，只恢复 $0.7\%$ 振幅内容——守恒可外挂，相互作用内容不可）。全部结论跨 $5$ 个随机
种子确认。我们同时诚实划界：该判别现象\textbf{不迁移到光滑连续系统}（Toda 格子为负结果），
故本文只主张离散可积。这不是把箱球系统包装成通用架构，而是证明“物理模型库”路线可以产生
可度量、抗作弊的机器学习样板。
\end{abstract}

% \tableofcontents

% ========================================================================
\section{引言}\label{sec:intro}
% ========================================================================

深度学习的架构大多是被\emph{搜索}出来、而非被\emph{选择}出来的。序列建模从 RNN 到 LSTM、
Transformer、状态空间模型，许多结构先在任务压力和工程试错中显出效果，理论解释与能力边界往往随后
才补；当一个主流范式撞到边界，我们手边通常并没有一座系统的候选结构库可查，造模型的常见做法仍是
在已有架构上改一点、把熟悉结构迁到新场景、或凭灵感攒一个设计再上 benchmark。本文所属的研究纲领
押一个不同的赌注：\textbf{物理学本身就是这样一座可以被系统开采的候选模型库}。物理积累了大量被反复
验证的成熟数学骨架——Ising／自旋玻璃、非平衡热力学／扩散、哈密顿／辛、张量网络、对称性等变、
重整化群、混沌、\textbf{可积系统}等——而 AI 已经不止一次从中拿到有效结构：扩散模型把非平衡热力学
搬进生成~\citep{sohldickstein2015deep, ho2020denoising}，Hopfield 网络与玻尔兹曼机把自旋玻璃与
Ising 带入联想记忆与学习~\citep{hopfield1982neural}，振子／哈密顿结构正在长成长程序列模型（coRNN、
LinOSS~\citep{rusch2021coupled, rusch2024oscillatory}）。过去这类成功多是零散采矿；我们想问的是能否
把它工程化——把物理根族逐个翻出来，做成可训练网络，放到适合它的任务上与主流模型硬碰，并诚实记录
它擅长与不擅长什么。

一篇“物理模型库”的宣言或综述并不足以说服机器学习读者，因此本文不试图一次建成整座库，而是选其中
一格走一遍完整的样板：从物理结构出发，做成可训练网络，设计一个能真正检验其结构优势的任务，用审计
对照排除硬编与外挂约束，并标出它的能力边界。我们开采的第一格是\textbf{可积／可逆系统}——它的意义
不在“用 AI 解物理题”，而在于它提供了一种很特殊的序列交互原型：系统可以发生强非线性交互，但交互
内容\emph{精确}保留（孤子彼此碰撞、穿过、相移，而振幅内容毫发无损）。与之相对，通用神经序列模型
能学到近似动力学，但在长程组合与长度外推下，一点点单步误差会累积，守恒量、可逆性乃至更丰富的结构
内容都会漂移。于是真正要检验的不是“把硬编的可积规则当模型、它能不能外推”（那是同义反复），而是：
\textbf{一个带可积归纳偏置、但关键规则仍须从数据中学出的神经模型，能否学到通用模型与外挂约束都得
不到的相互作用内容？}

我们的样板落在\textbf{有限容量箱球系统}上。普通箱球系统容易出现硬编泄露——一个无视进位的固定规则
即可解掉任务；把进位容量限到 $K$ 后规则变得非平凡：进位满载时球必须穿过，正确的放球决策必须用到
进位状态。我们也不止评估球数守恒，而评估\textbf{孤子振幅内容}，即孤子振幅的多重集；它比球数更深，
同样的球数可以对应不同的振幅分解（例如 $\{5,3,1\}$ 与 $\{4,4,1\}$ 都是 $9$ 个球），因此单把球数
守恒外挂到任意模型上并不能恢复它。在 $K=3$、组合 $T=8$ 的任务上，把这套可积 carrier 结构做进去的
模型在 $5$ 个随机种子上达到 $100\pm0\%$ 的逐位准确率、球数守恒率与振幅内容保持率；两个审计给出关键
证据：把可学习放球门换成写死规则后，振幅内容降到 $40\pm3\%$，说明这 $60$ 分的优势来自学习而非
硬编；把通用模型的输出投影到正确球数后，球数守恒升到 $100\%$、振幅内容却只有 $0.7\pm0.8\%$，说明
它不是标量守恒约束可以复制的（表~\ref{tab:tier1}）。我们同时诚实划界：这是一个\textbf{离散}可积
系统上的正面结果，而非完整的 benchmark 家族；在光滑连续的 Toda 格子上，同样的判别现象并不出现
（\S\ref{sec:boundary}）。

这里需要说清与既有工作的边界。“带左→右扫描的循环／状态追踪模型能长度外推、注意力不能”是成熟
结论~\citep{liu2023transformers, merrill2023parallelism, merrill2024illusion}，本文\emph{不}以“赢
Transformer”为贡献，Transformer 仅作无结构参照；用物理结构当归纳偏置也有先例（哈密顿
网~\citep{greydanus2019hamiltonian}、拉格朗日网~\citep{cranmer2020lagrangian}、可逆耦合
网~\citep{dinh2014nice, dinh2016density}），但它们多把物理\emph{当目标}或\emph{给保证}。本文的不同
在于测量可积结构是否让网络\emph{学}到一个通用模型与外挂约束都得不到的相互作用内容，并用泄露与外挂
两个审计把“真学 vs 硬编”“结构 vs 守恒”钉死。箱球系统本身是成熟的可积数学物理对象~\citep{takahashi1990soliton}，
但据我们所知，尚无人将其作为可学习的序列模型、以长度／组合外推与主流模型正面对照。

综上，本文的贡献有四点：(i)~把“物理模型库”这条路线落成一个端到端、可复现、抗作弊的机器学习样板，
而非停留在宣言；(ii)~在有限容量箱球系统上构造一个带可积 carrier 结构的可学习模型，并给出它在长程
组合下精确保持孤子振幅内容的证据（\S\ref{sec:main}）；(iii)~提出泄露审计与外挂对照两个方法论护栏，
分别把“真学 vs 硬编”“可积内容 vs 标量守恒”分开（\S\ref{sec:setup:audits}）；(iv)~诚实给出边界——
结果目前限于离散可积，连续 Toda 的负结果表明该机制不会自动迁移（\S\ref{sec:boundary}）。

% ========================================================================
\section{背景与设定}\label{sec:setup}
% ========================================================================

\subsection{箱球系统与有限容量变体}\label{sec:setup:bbs}
% 要点：BBS 是离散可积 CA；孤子=连续 1 块；有限容量 K 让规则非平凡（防 leak）。
箱球系统（Box-Ball System, BBS）是一维 $0/1$ 格点上的离散可积元胞自动机：一个“进位”
（carrier）从左向右扫描，拾取遇到的球、在空格处放下；一段连续的 $1$ 是一个\textbf{孤子}，
振幅 $=$ 块长，向右以正比于振幅的速度移动，大孤子追上小孤子后\emph{穿过}而各自振幅不变。
\textbf{有限容量箱球系统}把进位容量限制在 $K$：进位满（持有 $K$ 个球）时，遇到的球
\emph{穿过}而非被拾取。这一改动使单步规则\emph{非平凡}——正确的放球决策必须用到进位的当前
计数，因此一个“无视进位”的写死规则解不掉（这正是我们泄露审计的基础，见 \S\ref{sec:setup:audits}）。
\textbf{TODO：补一张 K 容量单步示意图或伪码。}

\subsection{相互作用内容：一个比球数更丰富的守恒量}\label{sec:setup:content}
% 要点：孤子振幅多重集是守恒不变量，且比标量球数丰富（{5,3,1} 与 {4,4,1} 同球数不同振幅）。
球总数是最表层的守恒量，是一个标量；它\emph{可以}被外挂到任意模型上（\S\ref{sec:main}）。
可积性的签名要丰富得多：\textbf{孤子振幅的多重集}
$\amp(x)=\{a_1\ge a_2\ge\cdots\}$ 在演化下\emph{精确不变}，而它\textbf{不能}由球总数恢复
——两个状态可以有相同的球数、不同的振幅多重集（如 $\{5,3,1\}$ 与 $\{4,4,1\}$ 同为 $9$ 球）。
我们把这个多重集作为衡量“是否保住了可积相互作用内容”的指标。它由把状态演化到孤子完全分离、
读取块长得到；我们数值核实了它在有限容量动力学下守恒、且系统可逆（$K=2/4/6$）。
% 成熟度：数值硬验证（自检通过）。

\subsection{两个审计：全文的诚实方法论}\label{sec:setup:audits}
% 护栏3 的方法论。leak 审计防“硬编冒充学习”，bolt-on 对照防“守恒冒充可积”。
\begin{Definition}[泄露审计 / carrier-blind]\label{def:leak}
把结构模型里\emph{可学习}的放球门换成一个\emph{写死}的规则（\texttt{emit}$=1-$cell，不训练），
其余结构不变。若写死版与学出版得分相当，则“学习”是空的、优势来自硬编；若写死版明显更差，
则优势是\emph{真学出来的}。
\end{Definition}

\begin{Definition}[外挂对照 / bolt-on]\label{def:bolton}
取一个通用模型，在推理时把它的输出\emph{投影}到“与输入球数相同”的约束上（保留概率最高的
$N$ 个格子，$N=$ 输入球数）。这把标量守恒量\emph{强行外挂}。若外挂后仍恢复不出相互作用内容，
则该内容\emph{外挂不出来}，是结构独有的。
\end{Definition}

% ========================================================================
\section{主结果：有限容量箱球上的可积外推（Tier 1）}\label{sec:main}
% ========================================================================
% 这是全文的肉。骨架：任务 → 排行表 → 两条读法（真学 60 分；深 0.7\%）。

\paragraph{任务。}
学习\emph{一}个有限容量箱球步，然后在多孤子状态上\textbf{组合}它、演化过多次碰撞
（训练用少孤子、短；测试用多孤子、长）。评测四个量：逐位准确率、球数守恒率、
孤子振幅内容的\emph{精确匹配率}与 \emph{IoU}。参赛者：结构模型（带结构守恒的进位扫描 +
可学习放球门，\S\ref{sec:setup:bbs}）；通用模型 GRU / LSTM / Transformer；外挂对照
（GRU $+$ 球数投影，定义~\ref{def:bolton}）；泄露审计（carrier-blind，定义~\ref{def:leak}）。

\begin{table}[h]
\centering
\caption{有限容量箱球系统上的可积外推排行榜（$K=3$，组合 $T=8$ 步，$5$ 个随机种子，均值 $\pm$
标准差）。结构模型在四个指标上均 $100\pm0$；泄露审计只到 $40\pm3$ 振幅内容（$60$ 分学习
margin，非泄露）；外挂对照钉住球数 $100\%$ 但只恢复 $0.7\%$ 振幅内容（守恒可外挂、内容不可）。}
\label{tab:tier1}
\begin{tabular}{lcccc}
\toprule
\textbf{参赛者} & \textbf{准确率\%} & \textbf{球数守恒\%} & \textbf{振幅内容-精确\%} & \textbf{振幅内容-IoU\%} \\
\midrule
\textbf{结构进位模型} & $\mathbf{100.0\pm0.0}$ & $\mathbf{100.0\pm0.0}$ & $\mathbf{100.0\pm0.0}$ & $\mathbf{100.0\pm0.0}$ \\
GRU（通用） & $91.7\pm1.2$ & $0.0\pm0.0$ & $0.0\pm0.0$ & $5.9\pm11.2$ \\
LSTM（通用） & $92.7\pm1.3$ & $1.7\pm3.3$ & $0.0\pm0.0$ & $13.5\pm11.6$ \\
Transformer（通用） & $71.6\pm10.9$ & $0.0\pm0.0$ & $0.0\pm0.0$ & $20.3\pm3.6$ \\
外挂对照（GRU $+$ 钉球数） & $89.2\pm2.5$ & $100.0\pm0.0$ & $0.7\pm0.8$ & $21.8\pm5.3$ \\
泄露审计（carrier-blind） & $88.9\pm0.5$ & $100.0\pm0.0$ & $40.0\pm3.2$ & $57.9\pm2.3$ \\
\bottomrule
\end{tabular}
\end{table}

\begin{Result}[真学，非硬编]\label{res:genuine}
% 成熟度：数值硬验证（5 种子）。这是实证发现，不是数学定理。
结构模型的优势是\emph{学出来的}：泄露审计（写死放球门）只达到 $40.0\pm3.2\%$ 振幅内容，
而学出的门达到 $100\%$——一个 $60$ 分的学习 margin。这不是平凡箱球上那种泄露（写死规则也
$100\%$）；这里写死规则明显更差，说明放球门确实学会了用进位状态。
\end{Result}

\begin{Result}[深，外挂不出来]\label{res:depth}
% 成熟度：数值硬验证（5 种子）。
把球数外挂到通用模型上，球数守恒率升到 $100\%$，但振幅内容只恢复 $0.7\pm0.8\%$。标量守恒量
可外挂，孤子振幅多重集不可；\emph{在本文测试的模型、训练预算与评测协议下}，只有带可积
carrier 结构的模型保得住，通用模型（未外挂）在两者上都崩（振幅内容 $0\%$）。
\end{Result}

% ========================================================================
\section{机理}\label{sec:mechanism}
% ========================================================================
% 要点：结构守恒让精确单步可学 → 组合不漂；bolt-on 恢复不出更丰富不变量；离散组合放大误差。
\textbf{TODO：正文。}三条：
(i)~结构守恒（进位记账）把单步规则约束到“保守”，使它可被学到\emph{精确}，从而组合多次不漂；
(ii)~通用模型只学到近似单步，离散组合把小误差放大，先丢内容、后丢球数；
(iii)~外挂只钉标量约束，无法钉住更丰富的振幅多重集（碰撞后的振幅配对）。

% ========================================================================
\section{支撑实验}\label{sec:support}
% ========================================================================
% 要点：box-ball exactness（+bolt-on 对照）、Margolus 迁移。都作支撑、不作主结果。

\subsection{箱球（有限容量）：精确性赢自由式，守恒可外挂}\label{sec:support:bbs}
% 成熟度：数值硬验证（多种子）。见 experiments/box_ball_learned_vs_transformer。
在长度外推的设定下（训练 $L=32$、测试到 $L=256$，$5$ 个种子），结构进位模型达到 $100\%$ 精确
外推，通用扫描模型（GRU / LSTM / SSM）随长度漂移；把守恒外挂到自由式模型上的对照能把守恒率
钉回 $100\%$，却\emph{补不回}准确率（表~\ref{tab:boxball}，差 $9$–$21$ 分）。这印证主结果的核心
区分：\textbf{守恒可外挂、精确性不可}。

\begin{table}[h]
\centering
\caption{有限容量箱球（$K=6$）长度外推，$L=256$（OOD）的逐位准确率 / 球数守恒率（$5$ 种子）。
结构模型 $100/100$；自由式随长度漂；外挂对照把守恒钉回 $100$ 但准确率仍停在 $91$（差 $9$ 分）。
Transformer 无扫描先验，仅作参照。详见 \texttt{experiments/box\_ball\_learned\_vs\_transformer}。}
\label{tab:boxball}
\begin{tabular}{lcc}
\toprule
\textbf{参赛者} & \textbf{准确率\%} & \textbf{球数守恒\%} \\
\midrule
\textbf{结构进位模型} & $\mathbf{100}$ & $\mathbf{100}$ \\
LSTM（自由） & $91$ & $6$ \\
GRU（自由） & $79$ & $0$ \\
Transformer（参照） & $64$ & $0$ \\
外挂对照（LSTM $+$ 钉守恒） & $91$ & $100$ \\
\bottomrule
\end{tabular}
\end{table}

\subsection{Margolus 块自动机：换个可逆 backbone 仍成立}\label{sec:support:margolus}
% 成熟度：数值硬验证（多种子 + bolt-on 对照）。见 experiments/margolus_block_ca。
同一套“结构守恒 $+$ 组合”配方迁移到一个结构不同的可逆系统（Margolus 块划分 CA，无长程进位；
步数随长度增长 $T\propto L$）：自由式模型即便单步学到 $\sim99\%$，长组合下守恒塌到 $\sim0$、准确率
掉到随机；把守恒外挂上去能钉回守恒率，却仍补不回准确率；结构模型全程 $100/100$
（表~\ref{tab:margolus}）。这说明主结果不是箱球/进位特有。

\begin{table}[h]
\centering
\caption{Margolus 块 CA（$5$ 种子）：单步准确率，以及组合到 $L=384$ 的准确率 / 守恒率。自由式
单步近乎学满，长组合下守恒崩、准确率掉到随机；外挂对照守恒钉回 $100$ 但准确率不恢复；结构块-CA
全程 $100/100$。详见 \texttt{experiments/margolus\_block\_ca}。}
\label{tab:margolus}
\begin{tabular}{lccc}
\toprule
\textbf{参赛者} & \textbf{单步 acc\%} & \textbf{组合 acc\% ($L{=}384$)} & \textbf{组合守恒\% ($L{=}384$)} \\
\midrule
\textbf{结构块-CA} & $\mathbf{100}$ & $\mathbf{100}$ & $\mathbf{100}$ \\
bi-LSTM（自由） & $98.7$ & $53$ & $1$ \\
Transformer（自由） & $99.7$ & $50$ & $0$ \\
外挂对照（bi-LSTM $+$ 钉守恒） & $98.7$ & $56$ & $100$ \\
\bottomrule
\end{tabular}
\end{table}

% ========================================================================
\section{边界：离散 $\to$ 连续不迁移（Toda 负结果）}\label{sec:boundary}
% ========================================================================
% 护栏2。诚实划界，不藏。见 experiments/toda_lattice。
% 成熟度：数值（负结果）。
让本文判别现象成立的机制\emph{依赖离散性}。在 Toda 格子（真·可积、但\textbf{光滑连续}）上，
一个自由式组合型步进器\emph{已经}把时间外推任务解得足够好（状态误差 $\sim0.003$），结构没有
可赢的 gap——光滑流的小步被自由式精确拿下、组合不累积，不存在离散箱球那种“单步差一点点、
$T\propto L$ 放大成灾难”的判别机制。我们诚实地把 Toda 记为负结果，并据此把本文主张\textbf{限定
在离散可积}。我们亦试过两条把可积结构\emph{学出来}的路（可学习作用-角模型、Neural-Lax），
均在连续设定下未能胜过外挂基线，一并作为负结果记录（附录~\ref{app:negatives}）。

% ========================================================================
\section{相关工作}\label{sec:related}
% ========================================================================
% 护栏1 的展开。三层：长度外推（撞的是能力结论）、物理归纳偏置（撞的是零件）、BBS（无人当模型）。
\textbf{TODO：三层展开，明确 delta（见护栏1）。}

% ========================================================================
\section{结语与展望}\label{sec:outlook}
% ========================================================================
% 要点：这是第一块砖；未来（更多系统/K 的家族、更深可积如带标签彩色内容）——明标为展望、暂缓。
本文给出一个在离散可积系统上、抗作弊（两个审计）、多种子确认的正面结果：可积结构让网络学到
通用模型与外挂都得不到的相互作用内容。这是把“物理可积模型作为可学习归纳偏置”落地的第一块砖，
而非完整的 benchmark 家族。\textbf{展望（暂缓，非本文主张）}：把该验证过的 Tier 1 扩成
跨多个离散可积/可逆系统、多容量的家族（含带颜色标签的彩色箱球，使不变量更丰富）；以及把“连续
可积（KdV / Lax）可学习化”的更深问题——目前尚为开放（Toda 负结果给了它一个诚实的下界）。

\appendix
% ========================================================================
\section{孤子振幅内容：提取协议与自检}\label{app:metric}
% ========================================================================
% 护栏3：主结果的指标必须写成固定 protocol、抗作弊，且对 true/pred 一视同仁。
振幅内容 $\amp(x)$ 由一个固定协议计算，对真值与模型预测\emph{同样}施加：

\begin{enumerate}
\item \textbf{演化到分离，读块长。} 把状态在右侧充分补零后按有限容量动力学演化，直到孤子完全
      分离——判据是块长多重集在再演化 $\Delta$ 步后\emph{不变}；此时读取各连续 $1$ 块的长度，
      得到排序多重集。分离所需步数按 $2|x|+60$ 起、每次翻倍直到稳定。
\item \textbf{补零量（安全界）。} 无界进位下最右球单步最多右移“总球数”格，故补零量取
      $(\text{步数})\times(\text{总球数})+|x|+10$，保证不越右边界。
\item \textbf{垃圾态处理。} 模型预测可能是非法态（球数错、非孤子结构）。为避免其“永不分离”导致
      的病态开销，翻倍最多 $4$ 次后返回当前最优估计。这对所有参赛者一致，不引入偏向。
\item \textbf{不偏向结构模型（关键）。} 该提取器是\emph{判别的}而非一律宽松：对垃圾输出给 $0\%$
      （通用模型），对部分正确给 $\sim40\%$（写死规则），对完全正确给 $100\%$（结构模型）。
      三档分明说明它度量的是“是否真的保住了孤子结构”，而非机械地偏袒某一方。
\end{enumerate}

\paragraph{自检覆盖。} 我们验证了 $\amp$ 在有限容量动力学下\emph{守恒}、且系统\emph{可逆}
（镜像法），覆盖容量 $K=2,4,6$、每档随机多孤子配置若干（含振幅超过 $K$、触发穿过的情形）；
教科书式的 $\{5,3,1\}$ 配置在演化下振幅内容恒为 $\{5,3,1\}$、球数恒为 $9$。
% 成熟度：数值硬验证（自检通过）。

% ========================================================================
\section{探路的负结果（诚实备忘）}\label{app:negatives}
% ========================================================================
% 护栏风格：把踩过的坑诚实收在附录，作为“为什么现在的设定是这一个”的证据。
\textbf{TODO：三段——}
(i)~\emph{平凡箱球 / 多孤子的泄露}：无容量限制时结构模型的放球门平凡（写死 $=1-$cell 即 $100\%$），
故换到有限容量以获得非平凡、真学的门；
(ii)~\emph{彩色箱球规则不可逆}：朴素彩色进位规则在自然镜像逆下只 $\sim86\%$ 可逆，非正则可积规则，
需晶体/组合-R 构造，故彩色作为未来 Tier 2；
(iii)~\emph{Toda 两发}（作用-角、Neural-Lax）：连续系统上自由式已解、结构无 gap。

% ========================================================================
% 参考文献
% ========================================================================
\bibliographystyle{unsrt}
\bibliography{references}

\end{document}
